% =====================================================================
% Manuscript for the Journal of Biomedical Informatics (Elsevier)
% Modeling Omics with Semantics for Dynamic Construction of
% Knowledge-Based Trans-Omic Networks
%
% Format: REVIEW mode (first submission)
%   - single column
%   - line numbers active (for reviewers)
%   - 11pt body, Times font
% Compile with: pdflatex -> bibtex -> pdflatex -> pdflatex
%               (or: latexmk -pdf manuscript_review.tex)
% =====================================================================
\documentclass[review,3p,times,11pt, ]{elsarticle} % twocolumn

% --- bibliography ----------------------------------------------------
% elsarticle handles natbib internally; pass numeric style.
\biboptions{numbers,sort&compress}

% --- math / tables / graphics ----------------------------------------
\usepackage{amsmath,amssymb,amsfonts}
\usepackage{booktabs}
\usepackage{tabularx}
\usepackage{multirow}
\usepackage{siunitx}
\usepackage{graphicx}
\usepackage{xcolor}
\usepackage{adjustbox}
\usepackage{rotating}
\usepackage{subcaption}

% --- code listings ---------------------------------------------------
\usepackage{listings}
\lstset{
    basicstyle=\ttfamily\scriptsize,
    keywordstyle=\color{black},
    commentstyle=\color{green!50!black},
    stringstyle=\color{red!70!black},
    numberstyle=\tiny\color{gray},
    stepnumber=1,
    numbersep=5pt,
    captionpos=b,
    breaklines=true,
    breakatwhitespace=false,
    showspaces=false,
    showstringspaces=false,
    showtabs=false,
    frame=single,
    tabsize=2
}

% --- algorithms (kept from original preamble) ------------------------
\usepackage{algorithm}
\usepackage{algpseudocode}
\usepackage{algorithmicx}

% --- misc utilities --------------------------------------------------
\usepackage{enumitem}

% --- hyperref (always last among packages) --------------------------
\usepackage{url}
\usepackage[hidelinks,colorlinks=false,unicode]{hyperref}

% --- show line numbers on the typeset PDF for the reviewers ----------
\usepackage{lineno}
\usepackage{comment}

% --- custom reference / url helpers ----------------------------------
\newcommand{\tabrefblue}[1]{\hyperref[#1]{Table~\ref*{#1}}}
\newcommand{\figrefblue}[1]{\hyperref[#1]{Figure~\ref*{#1}}}
\newcommand{\urlblue}[1]{\url{#1}}
% =====================================================================
\journal{Journal of Biomedical Informatics}

\begin{document}


\begin{frontmatter}

\title{MOGNN-TF: a heterogeneous graph neural network integrating transcription-factor regulation for pan-cancer molecular subtype classification}

%  **Proposta A — descrittiva e diretta (la più sicura)**
% > *MOGNN-TF: a heterogeneous graph neural network integrating transcription-factor regulation for pan-cancer molecular subtype classification*

% **Proposta B — il finding in primo piano**
% > *Transcription factors as a substitutable knowledge layer for multi-omics pan-cancer subtype classification with graph neural networks*

% **Proposta C — la formulazione "graph-first"**
% > *A three-layer heterogeneous graph of genes, miRNAs and transcription factors for pan-cancer subtype classification*


\author[unina]{Giovanni Maria De Filippis\corref{cor1}}
\ead{giovannimaria.defilippis@unina.it}

\author[unina]{Antonio M. Rinaldi}

\author[unina]{Cristian Tommasino}

\affiliation[unina]{%
  organization={Department of Electrical Engineering and Information Technology,
                University of Naples Federico II},
  addressline={Via Claudio 21},
  postcode={80125},
  city={Naples},
  country={Italy}}

\cortext[cor1]{Corresponding author.}

% =====================================================================
% ABSTRACT  —  v1, EN
% Sostituisce il blocco italiano + slot vuoti attualmente in bozza.
% Conta ~250 parole, distribuite ~50/80/80/40 sui quattro blocchi.
%
% Note di template:
%   - Se target è BMC Bioinformatics, sostituire i \textbf con \begin{abstract}
%     \backgroundtext{...} \resultstext{...} \conclusionstext{...} (bmcart).
%   - Per Bioinformatics (OUP): mantenere Motivation/Results/Availability/Contact,
%     spostando il blocco "Availability" nel suo slot apposito.
% =====================================================================
% --------------------------------------------------------------------
\begin{abstract}

Molecular subtype classification across the TCGA Pan-Cancer Atlas is
a difficult instance of supervised learning: $27$ heterogeneous classes,
strong imbalance, and roughly $10^4$ molecular features per sample
against $\sim$\,$9{,}000$ tumours with complete multi-omic profiles.
Graph neural networks (GNNs) inject prior biological knowledge into the
integration of mRNA, copy-number and miRNA data, but most existing models
restrict the prior to a single network type — typically protein--protein
interactions and, less often, miRNA--target predictions — model the
transcriptional regulatory layer only rarely, and are evaluated on binary
or few-class settings, leaving open whether a richer, multi-source prior
helps on the harder pan-cancer task. We present MOGNN-TF (Multi-Omics
Graph Neural Network with Transcription Factors), a heterogeneous GNN that
integrates three complementary sources of curated biological knowledge —
protein-level interactions from BioGRID, post-transcriptional regulation
from miRDB, and transcriptional regulation from TFLink — into a single
multi-relational graph where genes, miRNAs and transcription factors (TFs)
are first-class node types. On the 27-class iCluster pan-cancer taxonomy,
MOGNN-TF reaches a macro-F1 of $0.7918 \pm 0.0156$ on the test set,
outperforming two graph-based baselines (closest at $0.7703$), with the
improvement supported by Friedman and Holm-corrected Wilcoxon tests across
eleven seeds.
% --- Varianti alternative per la frase del confronto baseline (non scelte) ---
% Var. A (numeri espliciti, incl. baseline autoencoder/patient-similarity):
%   MOGNN-TF reaches a macro-F1 of $0.7918 \pm 0.0156$ on the test set,
%   against $0.7703$ for the closest published baseline and $0.6473$ for an
%   autoencoder-based patient-similarity baseline, with the improvement
%   supported by Friedman and Holm-corrected Wilcoxon tests across eleven seeds.
% Var. C (solo il numero MOGNN-TF, nessuna cifra delle baseline):
%   MOGNN-TF reaches a macro-F1 of $0.7918 \pm 0.0156$ on the test set,
%   significantly outperforming the competing baselines (Friedman and
%   Holm-corrected Wilcoxon tests across eleven seeds).
% ---------------------------------------------------------------------------- 
A four-axis ablation isolates the contribution of each prior
knowledge source and shows that TFs can substitute miRNAs in the graph with
no statistically significant loss in macro-F1, suggesting that the
transcriptional regulatory layer carries information partly redundant with
miRNA-target edges and could replace it when miRNA profiles are unavailable.

\end{abstract}

\begin{keyword}
Multi-omics integration \sep
Graph neural networks \sep
Cancer subtype classification \sep
Transcription factor regulation \sep
Pan-cancer \sep
TCGA Pan-Cancer Atlas \sep
Prior knowledge \sep
Heterogeneous graph
\end{keyword}

\end{frontmatter}

\linenumbers

% command fix (kept from previous version)
\renewcommand{\emph}[1]{\textit{#1}}
% 1. Rilassare i parametri di float (consigliato)
\renewcommand{\topfraction}{0.9}      % max frazione della pagina occupata da float in alto (default 0.7)
\renewcommand{\bottomfraction}{0.8}   % stesso per il bottom (default 0.3)
\renewcommand{\textfraction}{0.1}     % min frazione di testo richiesta (default 0.2)
\renewcommand{\floatpagefraction}{0.8}  % min riempimento per una "float page" (default 0.5)


% =====================================================================
% INTRODUCTION  —  v1, EN
% Sostituisce il blocco abstract italiano + le poche righe di introduzione 
% attualmente presenti. Da incollare DOPO \linenumbers e DOPO il
% \renewcommand{\emph}, prima della sezione "Related Work" o "Background".
% 

% Strategia per questa Introduction

% Struttura a funnel come da _paper_style_intructions_esteso_.md: problema → stato dell'arte e suo limite → gap → contributo → (outline minimale opzionale).
% Lunghezza target ~550 parole, cinque paragrafi asimmetrici.
% Citazioni selettive: 4 cluster, ~10 referenze totali. Lascio fuori la disamina dettagliata, che va in Related Work.
% Stile umanizzato: hedging (may, appear to, suggests), un paio di avversative, una frase corta per enfasi, una più lunga per argomentazione, voce attiva "we" nei punti d'azione.
% Niente vocabolario-spia AI (no delve, leverage, comprehensive, robust, pivotal, cutting-edge, paradigm).
% Posizionamento del gap reso falsificabile: "no prior work integrates X, Y, Z simultaneously on the 27-class iCluster taxonomy".


% =================================================================
\begin{comment}
    
**Note tecniche su cosa cambia rispetto alla bozza**

1. **Sostituisce** sia il blocco italiano dell'abstract (al momento contiene narrazione discorsiva non adatta) sia il primo paragrafo introduttivo. L'abstract va riscritto separatamente (TODO 4).
2. **Citazioni**: ho mantenuto tutte le chiavi BibTeX già presenti nella bibliografia. Le uniche due nuove sono:
 - `tonn_pancancer_grn_2025` — il preprint bioRxiv sul pan-cancer patient-specific GRN, da aggiungere alla `.bib` in TODO 3. Se preferisci non citare un preprint, posso sostituirla con una referenza più consolidata su TF e cancro (es. Bhagwat & Vakoc 2015, *Trends Cancer*, "Targeting transcription factors in cancer").
 - `liska_tflink_2022` — è la citazione di TFLink (Liska et al., *NAR* 2022). Verifico che sia in bib quando affrontiamo TODO 3; se manca, la aggiungo.
3. **Numero finale di MOGNN-TF**: ho riportato $0.7918 \pm 0.0156$ vs. $0.7703$ (Li24-GCN) e $0.6473$ (MoGCN). Il numero di MoGCN nell'abstract può sembrare crudele ma è quello che riportano i tuoi esperimenti — possiamo addolcirlo dicendo solo "outperforms two graph-based baselines" se preferisci.
4. **Outline finale**: l'ho tenuto, ma è il pezzo più "AI-like" della Introduction. Se l'editor del journal lo trova superfluo (BMC e Bioinformatics tendono ad accettarlo, alcune riviste no), lo togliamo in fase finale senza danni.

**Cose che ho deliberatamente NON fatto, e perché**

- *Non* ho scritto una "history of GNNs" né "since the advent of next-generation sequencing". Apertura diretta sul problema, come da style guide.
- *Non* ho elencato tutti i lavori GNN+multi-omics: la disamina cluster-by-cluster va in Related Work (TODO 2). Qui cito solo i rappresentanti necessari a definire il gap.
- *Non* ho usato connettivi-spia (*furthermore*, *moreover*, *additionally*, *in conclusion*) né verbi-spia (*delve*, *leverage*, *foster*, *harness*).
- *Non* ho fatto claim più forti dei risultati: "remain modest" è descrittivo onesto del Δ rispetto alla baseline modificata; "no statistically significant loss" è il fraseggio corretto del finding TF↔miRNA.

**Possibili variazioni che posso applicare al volo**

- Versione più assertiva (~450 parole) — taglio outline finale, accorcio §2.
- Versione più clinica — più peso al razionale clinico in apertura, meno al "label space size", utile se miriamo a JBI.
- Versione più formale — meno "frankly" e meno hedging, se preferisci un registro più rigido stile *Bioinformatics* main track.

Fammi sapere se va bene così o se vuoi una delle varianti, e poi proseguo con la Related Work (TODO 2) nel turno successivo. Se hai tempo, do un'occhiata anche tu alla lista di contributi: voglio essere sicuro che il punto 2 (sostituibilità TF↔miRNA) sia esattamente come lo vuoi inquadrare, perché è il finding più memorabile e i reviewer ci si fermeranno sopra.


\end{comment}



% =====================================================================
\section{Introduction}
\label{sec:introduction}

Classifying tumours by their molecular profile, rather than by their tissue of origin, is one of the more consequential shifts that pan-cancer genomics has produced over the last decade~\cite{weinstein2013tcga_pancancer,Hoadley2018CellOfOriginPanCan}. The taxonomy derived by Hoadley and colleagues from the TCGA Pan-Cancer Atlas separates roughly ten thousand tumours into 27 integrated iCluster subtypes that cut across the boundaries of histology and primary site, and that already differ substantially in prognosis and in the therapeutic decisions they suggest. The problem this raises for computational methods is not abstract: given mRNA, copy-number and miRNA profiles from a new sample, can we assign it to the correct integrated subtype with enough confidence to be useful downstream? The setting is, frankly, hard. The label space is large (27 classes, several with fewer than fifty samples), the feature space is on the order of $10^4$--$10^5$ molecular features per sample, and even the largest public cohort offers only $\sim$\,$9{,}000$ tumours with paired profiles across all three omic layers.

Deep multi-omics integration has been the dominant response, and graph-based methods in particular have gained ground. Approaches such as MOGONET~\cite{wang_mogonet_2021}, MoGCN~\cite{li_mogcn_2022}, SUPREME~\cite{kesimoglu_supreme_2022} and MOGAT~\cite{tanvir_mogat_2024} build patient-similarity graphs from one or more omic views and use a graph convolutional or attention-based model for downstream classification. A complementary line of work places the prior biological knowledge \emph{inside} the graph itself, by representing genes and miRNAs as nodes connected through curated interactions: Li and Nabavi~\cite{li_multimodal_2024} couple gene--gene interactions from BioGRID with miRNA--target relations, while Xiao et al.~\cite{xiao_graph_2023} and Tan et al.~\cite{tan_amogel_2025} explore the use of several prior knowledge networks in parallel. The intuition behind both lines is the same — message passing should respect known biology — but the gains over modality-naive baselines remain modest, and the question of \emph{which} prior is doing the work is rarely answered.

Two limitations are worth singling out. First, almost none of the prior-knowledge GNN frameworks for cancer subtyping model the \emph{transcriptional} layer explicitly. Transcription factors are central to gene regulation, and aberrant TF activity is implicated in roughly a third of human cancers~\cite{belova_pancancer_grn_2025}, yet TF--gene edges are usually absorbed into a generic protein--protein interaction backbone or ignored altogether. Second, evaluation tends to be carried out on binary or few-class settings (often a single tumour type, or BRCA-PAM50), which makes it hard to tell whether the apparent benefit of a more elaborate prior is real or a side-effect of an easier label space. Pan-cancer classification on the full 27-class iCluster taxonomy, with strong imbalance and partly overlapping subtypes, has received much less attention.

We address this gap by proposing \textbf{MOGNN-TF} (Multi-Omics Graph Neural Network with Transcription Factors). The model extends the multimodal GNN of Li and Nabavi~\cite{li_multimodal_2024} by adding an explicit transcriptional layer derived from TFLink~\cite{liska_tflink_2022}, so that the heterogeneous graph carries three biologically distinct sources of prior knowledge in parallel — protein-level interactions (BioGRID), post-transcriptional regulation (miRDB), and transcriptional regulation (TFLink) — and is evaluated, as a single integrated structure, on the full 27-class pan-cancer iCluster taxonomy. Our main contribution is this heterogeneous, multi-relational graph itself: a structure that jointly encodes gene--gene, miRNA--gene and TF--(gene/miRNA/TF) interactions, with genes, miRNAs and transcription factors as first-class node types, which to our knowledge has not been used before on pan-cancer subtype classification.

Around this graph, the paper does three further things. It runs a four-axis ablation (number of selected genes, sampling strategy, prior-knowledge composition, training fraction) that quantifies the contribution of each source separately; the most striking outcome is that transcription factors can replace miRNAs in the graph without a statistically significant loss in macro-F1, which suggests that the transcriptional layer carries information partly redundant with miRNA--target edges. It compares MOGNN-TF against two graph-based baselines~\cite{li_multimodal_2024,li_mogcn_2022} under Friedman and Holm-corrected Wilcoxon tests across eleven seeds, reaching a macro-F1 of $0.7918 \pm 0.0156$. And it releases a fully reproducible pipeline (Optuna-tuned hyperparameters, public data sources, Makefile-driven smoke and full reproduction paths) alongside the paper. Running through these analyses is a single explanatory thread, which we make explicit and test in Section~\ref{sec:discussion}: the benefit of a biological prior here appears to be mediated less by the choice of edge source than by how much connected structure survives feature selection, and the transcriptional layer helps precisely because it reconnects a graph that dimensionality reduction would otherwise fragment.

The rest of the paper is organised as follows: Section~\ref{sec:related_work} discusses related work, Section~\ref{sec:methods} describes data, graph construction and the architecture, Section~\ref{sec:results} reports the experimental evaluation and ablations, and Section~\ref{sec:discussion} discusses the findings and their limits.


% =====================================================================
% TODO 2 — versione precedente della lista dei contributi (rimossa su
% commento di Cristian: «rimuovi l'elenco, e metti in evidenza il primo
% punto come contributo principale dell'articolo. gli altri punti non
% sono un contributo»). Convertita in prosa qui sopra. Conservata per
% riferimento / eventuale reintroduzione come lista corta se il target
% primario diventa IEEE JBHI.
% ---------------------------------------------------------------------
% ...on the full 27-class pan-cancer iCluster taxonomy. The contributions are:
%
% \begin{itemize}\setlength{\itemsep}{1pt}
%     \item A heterogeneous graph that jointly encodes gene--gene, miRNA--gene and TF--(gene/miRNA/TF) interactions, which to our knowledge has not been used before on pan-cancer subtype classification.
%     \item A four-axis ablation study (number of selected genes, sampling strategy, prior-knowledge composition, training fraction) that quantifies the contribution of each source separately and shows that transcription factors can replace miRNAs in the graph without a statistically significant loss in macro-F1.
%     \item A robust statistical comparison against the published baseline of Li and Nabavi~\cite{li_multimodal_2024} and against MoGCN~\cite{li_mogcn_2022}, with Friedman and Holm-corrected Wilcoxon tests across eleven seeds; MOGNN-TF reaches macro-F1 $0.7918 \pm 0.0156$, against $0.7703$ and $0.6473$ for the two baselines.
%     \item A fully reproducible pipeline (Optuna-tuned hyperparameters, public data sources, Makefile-driven smoke and full reproduction paths) released alongside the paper.
% \end{itemize}
% =====================================================================


% ⚠️ Una nota, non un problema: la Discussion poggia molto sulla tesi che il guadagno è topologico (il TF layer riconnette il grafo frammentato dalla feature selection) — è il ponte concettuale tra R1 ed R2. La nuova Introduction non anticipa questo meccanismo (di proposito: in intro abbiamo solo accennato che "il layer trascrizionale porta informazione parzialmente ridondante"). Questo è coerente con una funnel-intro sobria, ma segnalo che il lettore incontra il concetto-chiave "connettività" solo in Discussion. Se in un secondo momento vuoi rafforzare il filo R1↔R2, l'aggancio andrebbe seminato in Intro — ma non è richiesto dal TODO 2 e l'hai voluta neutra, quindi lo lascio così. --> RISOLTO agginhengo "Running through these analyses is a single  ... fragment" alla fine del paragrafo di outline, per anticipare il filo conduttore che poi sviluppiamo in Discussion.
% =====================================================================
% RELATED WORK  —  v1, EN
% Sostituisce/integra la sezione attualmente abbozzata in italiano sui
% lavori correlati. Da inserire DOPO la sezione Introduction.
%
% Stile target: Bioinformatics (Oxford). Registro asciutto, voce attiva,
% cluster per approccio (non per paper), posizionamento esplicito alla fine.
% =====================================================================

\begin{comment}
Ho tutto. Procedo con la stesura della Related Work. Tre cluster:

Cluster A — patient-similarity / sample-graph approaches (il grafo è sui pazienti, non sulle entità biologiche).
Cluster B — biological-prior / feature-graph approaches (il grafo è sulle entità biologiche, gene/miRNA/proteina).
Cluster C — hierarchical e regulatory-informed (lavori che modellano regolazione, ma non in modo trans-omico integrato).

Più paragrafo finale di posizionamento esplicito di MOGNN-TF.

Mi sono permesso una nota di metodo: ho lasciato fuori dalla Related Work alcuni lavori meno centrali per non saturare (DeepMOCCA, MoSGAT, DGHNN, MoGCN-2022 non originale, MOLUNGN). Sono tutti citabili in Discussion o come riferimenti collaterali se servono.

## Cosa c'è di nuovo nella Related Work
Tre cluster per approccio, non per paper, come prescrivono le linee guida che mi hai dato. Ciascun cluster chiude con il limite strutturale, non con generiche "weakness". Il paragrafo finale di posizionamento è esplicito: triangolazione tra i tre cluster, non "our work is different".
Cosa ho deliberatamente fatto, con motivazione

HallmarkGraph è citato come "closest contemporaneous work" e distinto in due modi specifici: usa regulatory interactions ma è single-modality (solo mRNA), e modella una tassonomia gerarchica diversa (405 sottotipi su 26 cancer) — non è in concorrenza diretta con il setting 27-class iCluster. È onesto e prudente.
TF-DWGNet è chiamato per nome per disambiguare l'acronimo "TF". Sembra una sciocchezza ma per i reviewer di Bioinformatics è uno di quei dettagli che, se mancano, fanno pensare "non conoscono la letteratura recente". Una frase corta, "TF stands for Tensor Fusion, not transcription factor", risolve la cosa una volta per tutte.
MPKGNN (xiao_graph_2023) trattato come precedente più vicino al cluster B, e poi smontato sul punto specifico (manca la layer TF). Onesto e difendibile.
MoGCN è criticato sulla scalabilità richiamando il fatto che nei tuoi esperimenti hai dovuto ridurre le feature a varianza per farlo girare. Questo non è gratuito: è esattamente il tipo di osservazione che un reviewer apprezza perché si riferisce a un'evidenza concreta del tuo paper, non a un'opinione.

## Cosa NON ho fatto, con motivazione

Non ho citato survey come kang_roadmap_2022, baiao_technical_nodate, valous_graph_2024, wu2021comprehensive. Le survey vanno meglio in Background/Preliminaries che in Related Work — e Bioinformatics ha pochissimo spazio per survey citation. Le riservo per un eventuale paragrafo introduttivo della sezione Methods, se servono.
Non ho menzionato lavori solo tangenzialmente correlati (DeepMOCCA, MoSGAT, DGHNN, lavori su single-cell). Inflazionare la Related Work è il primo errore che un reviewer di Bioinformatics segna come "encyclopedic".
Non ho ripetuto numeri di performance dalla letteratura. Citarli senza poterli confrontare 1:1 col proprio setting confonde più di quanto chiarisca.

## Una raccomandazione operativa
Ti suggerisco di rileggere in particolare il sotto-paragrafo Patient-similarity GNNs (§3.1): è quello dove ho fatto la critica più affilata ("graph the model learns from is essentially a learned kernel"). Se preferisci un tono più morbido — alcuni reviewer di Bioinformatics vivono in quel mondo e potrebbero leggerla come polemica — posso ammorbidirlo in "patient-similarity graphs encode statistical proximity rather than mechanistic relationships", che dice la stessa cosa con meno mordente.

\end{comment}

\section{Related work}
\label{sec:related_work}

Work on graph-based multi-omics integration for cancer subtype classification
splits naturally into three families, distinguished by what the graph
represents and where the prior knowledge sits. We group the literature along
these lines, and close with a short statement of how MOGNN-TF differs from
each.

\subsection{Patient-similarity GNNs}
\label{ssec:rw_psn}

A first family treats each patient as a graph node and connects nodes by
similarity across one or more omic views. MOGCN~\cite{li_mogcn_2022} couples
an autoencoder with a Similarity Network Fusion step to build the patient
graph, while SUPREME~\cite{kesimoglu_supreme_2022} learns per-modality
patient embeddings and merges them downstream; MOGAT~\cite{tanvir_mogat_2024}
adopts the same backbone with attention, and
DeepMoIC~\cite{wu_deepmoic_2024} uses a deeper GCN with residual connections
to stabilise training over the resulting similarity network. The recent
MO-GCAN~\cite{dou_mo-gcan_2025} layers an attention mechanism on top of the
fused patient graph and reports competitive results on BRCA and KIPAN. These
methods are attractive because they sidestep the need for a curated
biological network, but the graph they learn from is essentially a learned
kernel: edges encode statistical proximity between patients, not mechanistic
relationships between molecular entities, and the model is not exposed to
prior biological structure during message passing. 
% --- SOFTER ALTERNATIVE (uncomment to use, comment out the version above) ---
% but the resulting graph encodes statistical proximity between patients
% rather than mechanistic relationships between molecular entities, which
% means the model is not directly exposed to prior biological structure
% during message passing.
% ---------------------------------------------------------------------------
Patient-similarity methods also tend to scale poorly with the number of samples, since the
similarity matrix is dense; in pan-cancer settings with $\sim$\,$10^4$
tumours, even MOGCN had to be evaluated on a variance-pruned subset
in our experiments. The interpretation problem is harder, too: nodes are
patients, so a saliency analysis returns ``which patients matter'', not
``which genes or regulators matter''.

\subsection{Biologically-informed feature graphs}
\label{ssec:rw_biograph}

A second family takes the opposite view and puts genes (and, where
available, miRNAs) as nodes, drawing edges from curated interaction
databases. MOGONET~\cite{wang_mogonet_2021} is the canonical reference
here, with a per-omic GCN whose graphs come from feature correlations but
whose later fusion exploits view-correlation discovery. Li and
Nabavi~\cite{li_multimodal_2024} —- the baseline we extend -— take a
sharper line and inject a biological gene--gene interaction network
(BioGRID) and miRNA--target edges directly into a heterogeneous graph,
attaching mRNA, CNV and miRNA features to the corresponding nodes; we
borrow the multimodal scaffolding from this work. AMOGEL~\cite{tan_amogel_2025}
combines a curated PPI/KEGG backbone with edges mined by Association Rule
Mining and a GAT classifier, evaluated on BRCA and KIPAN subtypes.
MPKGNN~\cite{xiao_graph_2023} pushes the idea further by stacking
\emph{multiple} prior knowledge graphs in parallel and learning a soft
combination, which is conceptually the closest precedent to our work.
Two things, however, limit this line. First, the prior is usually
restricted to PPI and miRNA--target edges: transcriptional regulation,
which is arguably the most direct mechanism behind subtype-defining
expression programs, is either folded into the PPI backbone or absent.
Second, evaluation is almost always carried out on binary or
few-class problems on a single tumour type (BRCA-PAM50 dominates the
benchmarks), so the question of whether a richer prior helps when the
label space is large and imbalanced -— the pan-cancer setting -— is left
open. A recent comparative study~\cite{alharbi_comparative_2025} confirms
the observation, noting that gains from heavier biological priors tend to
shrink as the classification task becomes harder.

\subsection{Regulatory- and hallmark-informed models}
\label{ssec:rw_regulatory}

A third, smaller family bakes transcriptional or pathway-level structure
into the model. HallmarkGraph~\cite{zhang_hallmarkgraph_2025} (the closest
contemporaneous work) uses gene regulatory interactions and the ten cancer
hallmark axes as an inductive bias to classify a 26-cancer, 405-subtype
hierarchical taxonomy from transcriptome data alone, reporting strong
results but restricting itself to mRNA and to single-modality input.
Belova et al.~\cite{belova_pancancer_grn_2025} go in a complementary
direction, reconstructing patient-specific GRNs across the TCGA pan-cancer
cohort and using them to stratify tumours by regulatory state; the
analysis is unsupervised and not aimed at the iCluster taxonomy, but it
documents how much subtype-relevant variability sits at the TF level.
Outside cancer subtype classification, TF-aware GNNs have appeared in
related tasks such as cancer-gene prediction~\cite{xiao_graph_2023}
and driver discovery, where TF--target edges are used as one of several
priors but without a controlled comparison against alternatives. We note,
incidentally, that ``TF'' in
TF-DWGNet~\cite{yang_tfdwgnet_2026} stands for \emph{Tensor Fusion},
not transcription factor; the model learns its graph in a supervised
tree-based fashion and does not use regulatory priors. The distinction
matters because, at the time of writing, no published GNN framework for
pan-cancer molecular subtyping integrates TF--target regulation as an
explicit, ablation-controlled prior alongside protein-level and
post-transcriptional knowledge.

\subsection{Position of MOGNN-TF}
\label{ssec:rw_position}

MOGNN-TF sits at the intersection of these three lines. We retain the
biologically-informed feature graph of the second family (genes, miRNAs
and now TFs as nodes; curated edges from BioGRID, miRDB and TFLink),
borrow the multimodal scaffolding of Li and Nabavi~\cite{li_multimodal_2024},
and apply it to the 27-class iCluster pan-cancer taxonomy of
Hoadley et al.~\cite{Hoadley2018CellOfOriginPanCan} -- the same hard,
imbalanced label space that most prior-knowledge GNNs avoid. Unlike
MPKGNN~\cite{xiao_graph_2023}, which combines several priors but does not
include TF regulation, we add transcription factors as a first-class
node type with their own edges (TF--gene, TF--miRNA, TF--TF), and unlike
HallmarkGraph~\cite{zhang_hallmarkgraph_2025} we use the regulatory layer
inside a true multi-omic graph (mRNA + CNV + miRNA), not as a constraint
on a transcriptome-only model. The four-axis ablation we report
(Section~\ref{sec:results}) is designed to make the contribution of
\emph{each} knowledge source separable -- a question that, to our
knowledge, the existing literature has not answered for this setting.
% =====================================================================
% MATERIALS AND METHODS — v1, EN
% Parte 1 di 2: §2.1 Data + §2.2 Preprocessing
% Da inserire DOPO la sezione Related Work.
% La parte 2 (Graph construction + Architecture + Training) arriva nel
% turno successivo.
% =====================================================================

\begin{comment}



## Tre cose da verificare quando hai un momento

- [RISOLTO] Numero N_1 simbolico. Sostituito con il valore reale verificato dai dati: 17.944 geni (intersezione CNV∩Expression dopo armonizzazione ID).
- [RISOLTO] Median imputation. Verificato sul codice: NON esiste median imputation; la pipeline usa solo fillna(0) per tutte le omiche. Frase rimossa e paragrafo "Missing values" riscritto come zero-fill. Corretta anche l'affermazione errata "CNV filled with -1" (il codice usa 0 anche per CNV).
- Sample filter 01/11/06. Ho dichiarato che usate solo 01 (primary solid tumour) di default. Verifica che questo sia coerente col config_final.yml del repo, perché è esattamente il tipo di dettaglio che i reviewer vanno a controllare.
\end{comment}



\section{Materials and methods}
\label{sec:methods}

MOGNN-TF is an end-to-end framework that turns three multi-omic views
and three curated interaction networks into a single subtype
prediction. The pipeline, shown in Figure~\ref{fig:framework}, proceeds
in five stages. We start from gene expression, copy number and miRNA
profiles of the TCGA Pan-Cancer cohort, together with the iCluster
labels. These matrices are preprocessed and harmonised to a common
gene identifier space, and their feature axis is aligned to the
prior-knowledge networks, so that the molecular features that survive
selection are exactly those that can be placed on the graph. The
aligned features then become the nodes of a single heterogeneous,
multi-relational graph that encodes the three sources in parallel —
protein-level interactions from BioGRID, post-transcriptional
regulation from miRDB, and transcriptional regulation from TFLink —
with genes, miRNAs and transcription factors as distinct node types. A
graph neural network encodes this structure, and a multilayer
perceptron maps the resulting representation to the 27 iCluster
subtypes.

The model itself (Section~\ref{ssec:architecture}) has four jointly
trained parts: a dimension-increase projector that lifts the
heterogeneous node features into a shared latent space, a relational
GNN encoder that propagates information across the graph, a
reconstruction decoder that regularises the latent representation, and
a shallow parallel branch that bypasses the graph; their outputs are
fused before the classifier. The remaining subsections follow the flow
of the data, reading the framework of Figure~\ref{fig:framework} from
left to right: Section~\ref{ssec:data} describes the omic and
prior-knowledge sources, Section~\ref{ssec:preprocessing} the
preprocessing and prior alignment, Section~\ref{ssec:graph} the
construction of the heterogeneous graph, Section~\ref{ssec:architecture}
the network architecture, and Section~\ref{ssec:training} the training
and evaluation protocol.

% TODO[fig]: nel PNG, rinominare il titolo del blocco centrale
%   "Li24 Architecture" -> "MOGNN-TF" (ora questa figura apre Methods
%   come roadmap del framework, non come architettura del baseline).
\begin{figure}[t]
    \centering
    \includegraphics[width=\linewidth]{images/framework_buono_finale.png}
    \caption{Overview of the MOGNN-TF framework, read left to right.
    Raw multi-omics data (gene expression, copy number, miRNA) and the
    three prior-knowledge interaction networks (BioGRID, miRDB, TFLink)
    are preprocessed and feature-selected, then assembled into a single
    heterogeneous, multi-relational graph in which genes, miRNAs and
    transcription factors are distinct node types. A dimension-increase
    projector lifts the node features into a shared latent space, a
    heterogeneous GNN encoder (GCN or GAT) learns subtype-relevant
    representations, and a reconstruction decoder together with a shallow
    parallel branch preserve global omic patterns; the fused
    representation feeds an MLP that outputs the molecular subtype. 
    % The per-component details are given in Sections~\ref{ssec:graph}
    % and~\ref{ssec:architecture}. 
    Modules introduced in this work are
    shown in green, those modified from the baseline of Li and
    Nabavi~\cite{li_multimodal_2024} in yellow, and the unchanged
    baseline modules in blue.}
    \label{fig:framework}
\end{figure}

% ---------------------------------------------------------------------
\subsection{Data}
\label{ssec:data}

We use the multi-omics resources of the TCGA Pan-Cancer
Atlas~\cite{weinstein2013tcga_pancancer} as harmonised by the UCSC Xena
platform~\cite{goldman2020xena}. Three molecular views are retained:
gene expression (RNA-seq, EB++ batch-corrected), gene-level copy number
variation from GISTIC2, and mature-strand miRNA expression (also EB++
batch-corrected). Both expression matrices are provided as
$\log_2(x{+}1)$ values; CNV is left on the GISTIC2 integer scale. After
basic format harmonisation, the raw matrices contain $11{,}069$ samples
$\times$ $20{,}530$ genes for mRNA, $10{,}845 \times 24{,}776$ for CNV,
and $10{,}824 \times 743$ for miRNA.

Class labels come from the 27-class iCluster pan-cancer taxonomy of
Hoadley et al.~\cite{Hoadley2018CellOfOriginPanCan}, downloaded from the
same Xena hub. The label file covers $9{,}666$ samples assigned to one
of clusters C1--C28 (C24 absent); the resulting label distribution is
strongly imbalanced, with the largest cluster containing roughly
$10\times$ more samples than the smallest, and four clusters with fewer
than fifty samples each. Three sources of prior knowledge are then
ingested: gene--gene protein interactions from BioGRID release
5.0.250~\cite{stark_biogrid_2006}, miRNA--target predictions from
miRDB~\cite{chen_mirdb_2020}, and human TF--target regulatory edges from
TFLink~\cite{liska_tflink_2022} (simple format, high-confidence subset).
The exact dataset URLs and versions are listed in the project repository.

% ---------------------------------------------------------------------
\subsection{Preprocessing}
\label{ssec:preprocessing}

The preprocessing pipeline, summarised in
Figure~\ref{fig:preprocessing}, has four objectives: align samples
across modalities, harmonise gene identifiers, handle missing values,
and align the resulting feature space to the prior-knowledge networks.

\paragraph{Sample alignment.}
The three omic matrices are first transposed so that rows correspond to
samples and columns to molecular features, and intersected on the
common sample set; samples without an iCluster label are discarded. By
default we retain only primary solid tumours (TCGA barcode prefix
\texttt{01}); the pipeline exposes a configurable filter to also include
\texttt{11} (solid tissue normal) and \texttt{06} (metastatic), which we
do not use in the experiments reported here.

\paragraph{Gene identifier harmonisation.}
Expression and CNV use different identifier conventions
(ENSEMBL IDs in expression, mixed annotations in CNV). We map both to
official HGNC symbols, remove duplicate rows, and convert legacy aliases
where needed. After harmonisation, mRNA and CNV are restricted to the
intersection of their gene sets, producing two matrices on a common
gene axis of $17{,}944$ genes. The miRNA matrix is left at $743$
features since it shares no identifier space with the gene matrices.

\paragraph{Missing values.}
We drop genes with more than $30\%$ missing values across samples; the
remaining gaps are zero-filled, following common practice on TCGA
data~\cite{zhang_mosgraphgen}. For CNV, zero coincides with the GISTIC2
neutral level (``no detected alteration''), and for the
$\log_2(x{+}1)$ expression values it corresponds to the no-expression
baseline, so the imputed value is biologically conservative in both
omics.

\paragraph{Prior-knowledge alignment.}
The last step restricts the omic matrices to entities that also appear
in the prior-knowledge networks. Genes are intersected with the union
of BioGRID and TFLink gene vocabularies; miRNAs are intersected with
the miRDB target catalogue, after standardising the identifier format
(e.g.\ \texttt{hsalet7a23p} $\to$ \texttt{hsa-let-7a-3p}). After this
step the working matrices contain $15{,}942$ genes for both mRNA and
CNV, and $688$ miRNAs, on the same $M$-sample axis. These are the
inputs to the downstream variance-based feature selection and graph
construction (Section~\ref{ssec:graph}), which reduce the working
dimension further to the operating set of $700$ genes, $100$ miRNAs and
$200$ TFs used in the reported experiments.

\begin{figure}[!t]
    \centering
    \includegraphics[width=\linewidth]{images/preprocessing_v4.png}
    \caption{Preprocessing pipeline. Raw multi-omics matrices from
    UCSC Xena are sample-aligned, gene-id harmonised, imputed and
    finally aligned to BioGRID, miRDB and TFLink, producing the working
    matrices used for graph construction.
    % --- TODO: redraw with English labels (currently in Italian) ---
    }
    \label{fig:preprocessing}
\end{figure}



% =====================================================================
% MATERIALS AND METHODS — v1, EN
% Parte 2 di 2: §2.3 Graph construction + §2.4 Architecture + §2.5 Training
% Da inserire DOPO §2.2 Preprocessing (TODO 5).
% =====================================================================

% ---------------------------------------------------------------------
\subsection{Heterogeneous graph construction}
\label{ssec:graph}

After preprocessing, the same heterogeneous graph $G = (V, E)$ is shared
across all samples; only the node features change. We define three node
sets corresponding to the three biological entities used by the model:
$V_G$ for genes, $V_M$ for miRNAs, and $V_T$ for transcription factors,
with $|V_G| = 700$, $|V_M| = 100$ and $|V_T| = 200$ in the final
configuration (Section~\ref{ssec:training}). Genes and TFs are kept as
disjoint node sets even though TFs are themselves genes: this allows the
model to treat regulatory and protein-level interactions as semantically
distinct edge types, and is the same convention adopted in
TFLink~\cite{liska_tflink_2022}.

Conceptually, the graph is organised in three biological layers
(Figure~\ref{fig:graph_concept}): a gene layer connected by
protein-level interactions (BioGRID), a miRNA layer linked to its target
genes (miRDB), and a transcription-factor layer connected to the genes
and miRNAs it regulates (TFLink). Intra-layer edges capture interactions
between entities of the same type, while inter-layer edges encode the
post-transcriptional and transcriptional regulation that ties the three
levels together; it is this third, transcriptional layer that
distinguishes MOGNN-TF from the two-level gene--miRNA graph of Li and
Nabavi~\cite{li_multimodal_2024}. The remainder of this section
formalises the node sets, the edges and the assembled adjacency that
realise this layered structure.

\begin{figure}[t]
    \centering
    \includegraphics[width=0.85\linewidth]{images/graph_mo-1.png}
    \caption{Conceptual structure of the heterogeneous graph. The three
    biological layers — genes connected by gene--gene interactions (GGI,
    BioGRID), miRNAs by miRNA--target relations (miRDB), and transcription
    factors by TF--target regulation (TFLink) — are shown with
    representative entities. Solid lines are intra-layer edges; dashed
    lines are the inter-layer regulatory edges that link the levels. 
    %Node types are distinguished by colour. This is a schematic; the dense     graphs that survive feature selection are shown in Figure~\ref{fig:graphs}.
    }
    \label{fig:graph_concept}
\end{figure}
% TODO[fig]: rigenerare graph_mo-1.png prima della submission:
%   (1) color-safety (TODO 7): i nodi usano rosso (gene) + verde (miRNA),
%       combinazione NON color-blind-safe -> ricolorare (es. blu/arancione/
%       magenta) o aggiungere forme/etichette ai tipi di nodo;
%   (2) refuso label: "Transcription Factoris Level" -> "Transcription
%       Factor Level".

\paragraph{Feature selection.}
Variance-based feature selection is applied separately to each node
type, following common practice in this
literature~\cite{li_multimodal_2024,xiao_graph_2023}. For genes, we rank
all $15{,}942$ aligned genes by their expression variance on the
training set and retain the top $N_g$. The same procedure is repeated
on miRNAs (top $N_m$ by miRNA-expression variance) and on TFs (top
$N_t$ by expression variance, restricted to the TFLink vocabulary). The
choice $N_g = 700$, $N_m = 100$, $N_t = 200$ is the operating point
from the Optuna study (Section~\ref{ssec:training}); the ablation on
$N_g$ is reported in Section~\ref{sec:results}. One subtle effect is
worth flagging here: because feature selection is done independently on
genes and TFs, some entities appear in both sets. For the configurations
we use this overlap is $47$ ($N_g = 700$, $N_t = 200$) and grows
roughly linearly with $N_g$. We keep the duplicates as distinct nodes,
because their roles in the graph (target vs.\ regulator) and their edge
neighbourhoods are different, but we acknowledge that this is a design
choice; collapsing them is a sensible alternative that we did not
explore.

\paragraph{Edges.}
Edges come from three curated sources. Gene--gene edges $E_{GG}$ are
the BioGRID release 5.0.250 human PPI network~\cite{stark_biogrid_2006},
restricted to the gene subgraph induced by $V_G \cup V_T$.
miRNA--target edges $E_{MG}$ come from miRDB~\cite{chen_mirdb_2020};
we keep predictions with score $\geq 80$ to retain only the
high-confidence subset. TF--target edges $E_{TG}$ come from
TFLink~\cite{liska_tflink_2022} (human, simple format,
small-scale-only filter). All edges are treated as undirected and
binary; we do not use edge weights, partly to keep the formulation
consistent with Li and Nabavi~\cite{li_multimodal_2024}, partly because
the source databases use heterogeneous confidence scales that would not
combine cleanly. Block-level adjacencies for the remaining relations
are derived from the explicit ones:
%
\begin{equation}
    \mathbf{A}_{MM} = \mathbb{I}\!\left((\mathbf{A}_{GM})^{\!T}\mathbf{A}_{GM} > 0\right),
    \label{eq:Amm}
\end{equation}
\begin{equation}
    \mathbf{A}_{TT} = \mathbb{I}\!\left(\mathbf{A}_{TG}\mathbf{A}_{TG}^{\!T} > 0\right),
    \label{eq:Att}
\end{equation}
\begin{equation}
    \mathbf{A}_{TM} = \mathbb{I}\!\left(\mathbf{A}_{TG}\mathbf{A}_{GM} > 0\right),
    \qquad
    \mathbf{A}_{MT} = \mathbf{A}_{TM}^{\!T}.
    \label{eq:Atm}
\end{equation}
%
The intuition is straightforward: two miRNAs are linked if they share
at least one target gene; two TFs are linked if they share at least one
target; a TF is linked to a miRNA if the TF regulates a gene that the
miRNA also targets. When a block is excluded by an ablation, intra-omic
blocks are replaced by the identity matrix and inter-omic blocks by the
zero matrix. The full adjacency is then assembled as
%
\begin{equation}
    \mathbf{A} =
    \begin{bmatrix}
        \mathbf{A}_{GG} & \mathbf{A}_{GM} & \mathbf{A}_{GT} \\
        \mathbf{A}_{MG} & \mathbf{A}_{MM} & \mathbf{A}_{MT} \\
        \mathbf{A}_{TG} & \mathbf{A}_{TM} & \mathbf{A}_{TT}
    \end{bmatrix},
\end{equation}
%
extending the two-level (gene + miRNA) formulation of Li and
Nabavi~\cite{li_multimodal_2024} to the three-level case that
incorporates transcriptional regulation explicitly.

\paragraph{Node features.}
The graph topology is shared across samples; the node-feature matrix
$\mathbf{X}^{(s)} \in \mathbb{R}^{(N_g + N_m + N_t) \times 2}$ is
sample-specific. For each tumour $s$, gene and TF nodes carry a
$2$-dimensional feature vector $[\textrm{expr}_i^{(s)},
\textrm{cnv}_i^{(s)}]$ (expression and copy number); miRNA nodes carry
the miRNA expression value, padded with a zero to match the shared
feature dimension. Per-omic min--max normalisation is applied separately
on each split, with parameters estimated on the training fold only to
avoid leakage. The resulting per-sample graph, together with the
post-selection gene--TF, gene--miRNA and full graphs, is shown in
Figure~\ref{fig:graphs}.

\begin{figure}[H]%[!t]
    \centering
    \begin{subfigure}{0.48\linewidth}
        \centering
        \includegraphics[width=\linewidth]{images/gene_mirna_graph.png}
        \caption{Gene--miRNA subgraph.}
        \label{fig:graph_gm}
    \end{subfigure}\hfill
    \begin{subfigure}{0.48\linewidth}
        \centering
        \includegraphics[width=\linewidth]{images/gene_tf_graph.png}
        \caption{Gene--TF subgraph.}
        \label{fig:graph_gt}
    \end{subfigure}\\[6pt]
    \begin{subfigure}{0.55\linewidth}
        \centering
        \includegraphics[width=\linewidth]{images/full_graph.png}
        \caption{Full gene--miRNA--TF graph.}
        \label{fig:graph_full}
    \end{subfigure}
    \caption{Heterogeneous graph used by MOGNN-TF after feature selection
    ($N_g{=}700$ genes, $N_t{=}200$ TFs, $N_m{=}100$ miRNAs).
    Genes are shown in red, miRNAs in green and TFs in blue.
    Adding the TF layer (panel c) substantially improves the
    connectivity of the subgraph that survives feature selection.}
    \label{fig:graphs}
\end{figure}

% ---------------------------------------------------------------------
\subsection{Architecture}
\label{ssec:architecture}

We now detail the network sketched in the right-hand portion of
Figure~\ref{fig:framework}. The architecture
builds on the multimodal scaffolding of Li and
Nabavi~\cite{li_multimodal_2024} and extends it to handle the new TF
node type. The model has four parts that are trained jointly: a
\emph{dimension-increase} projector that lifts the heterogeneous node
features into a shared latent space, a \emph{GNN encoder} that
propagates information across the heterogeneous graph, a
\emph{decoder} that reconstructs the omics input from the encoded
representation, and a \emph{shallow parallel branch} that bypasses the
graph entirely. The four outputs are concatenated by an
\emph{embedding-fusion} block and fed to a multilayer perceptron (MLP)
classifier that produces the subtype prediction.

\paragraph{Dimension-increase projector.}
The three node types carry input vectors of different size and
semantics: $\mathbf{X}_G \in \mathbb{R}^{N_g \times 2}$ (expression +
CNV), $\mathbf{X}_T \in \mathbb{R}^{N_t \times 2}$ (expression + CNV,
since TFs are themselves genes), and $\mathbf{X}_M \in \mathbb{R}^{N_m
\times 1}$ (miRNA expression). Following~\cite{li_multimodal_2024}, a
type-specific stack of fully-connected layers (linear $\rightarrow$
ReLU blocks, hidden size $\theta_2$) projects each block into a
common latent dimension; we add a third branch for TF nodes that
mirrors the gene branch. The three projected blocks are then stacked
back into a single node-feature matrix of shape
$(N_g + N_m + N_t) \times \theta_2$.

\paragraph{GNN encoder.}
Two graph convolution layers are stacked on top of the projected
features, separated by LeakyReLU activations and dropout. We support
both GCN~\cite{kipf_semi-supervised_2017} and
GAT~\cite{velickovic_graph_2018} convolutions, chosen at configuration
time; the comparison between the two is one of the questions we
investigate in the ablations (Section~\ref{sec:results}). The
graph-level representation is obtained by global \texttt{max} pooling
over all nodes, producing a single vector $\boldsymbol{\theta}_1$ of
size equal to the GNN hidden dimension.

\paragraph{Decoder and shallow branch.}
$\boldsymbol{\theta}_1$ feeds a symmetric decoder that reconstructs the
per-node omics features; this acts as an autoencoder-like
regulariser, encouraging the latent representation to retain biological
content beyond what is needed for classification alone. In parallel, a
shallow MLP (the ``parallel branch'' of~\cite{li_multimodal_2024})
takes the concatenated raw omics as input and produces a second
representation $\boldsymbol{\theta}_2$ that is not exposed to the
graph. The two vectors $\boldsymbol{\theta}_1, \boldsymbol{\theta}_2$
are concatenated and passed to a two-layer MLP classifier with
softmax output over the 27 subtype classes.

\paragraph{Loss.}
We train the whole model jointly with a composite loss
%
\begin{equation}
    \mathcal{L} = \mathcal{L}_{\text{CE}}(\hat{y}, y)
                + \lambda \, \mathcal{L}_{\text{rec}}(\hat{\mathbf{X}}, \mathbf{X})
                + \mu \, \|\boldsymbol{\theta}_{\text{GNN}}\|_2^2,
\end{equation}
%
where $\mathcal{L}_{\text{CE}}$ is the (unweighted) cross-entropy on
the subtype labels and $\mathcal{L}_{\text{rec}}$ is the reconstruction
MSE on the node features. The third term is an explicit L2 penalty on
the GNN parameters with coefficient $\mu = 5 \times 10^{-4}$, applied
in addition to the AdamW weight decay because the two terms target
different parameter subsets in our implementation. Class imbalance is
handled at the sampler level rather than through the loss
(Section~\ref{ssec:training}); the weight $\lambda$ is fixed across
the reported experiments at the value selected by the Optuna search.

% ---------------------------------------------------------------------
\subsection{Training and evaluation protocol}
\label{ssec:training}

\paragraph{Data splits.}
After preprocessing, the working cohort contains $9{,}355$ tumours
with complete mRNA, CNV and miRNA profiles. We split this cohort
stratified by subtype into $60\%/20\%/20\%$ for train, validation and
test, freeze the split indices on disk, and reuse them across all
experiments. The split seed is fixed at \texttt{42}, so that all
comparisons across runs operate on identical train/validation/test
partitions; eleven additional seeds (independent from the split seed)
control model initialisation, batch order and the weighted sampler.
The variability across these eleven seeds is what feeds the
Friedman/Wilcoxon analyses in Section~\ref{sec:results}.

\paragraph{Optimisation and early stopping.}
We use AdamW with learning rate $8 \times 10^{-4}$, weight decay
$1 \times 10^{-4}$, batch size $64$ and dropout $0.2$, with a cosine
annealing learning-rate schedule. Training is capped at $100$ epochs
but stops earlier through early stopping on validation macro-F1, with
patience $10$. When a new best validation score is reached, both model
weights and optimiser/scheduler state are checkpointed to disk to
allow exact resumption. Table~\ref{tab:hyperparams} summarises the
final configuration.

\begin{table}[!t]
\centering
\caption{Final MOGNN-TF configuration (Optuna trial 18, selected by
validation macro-F1). The full search grid is reported in the project
repository.}
\label{tab:hyperparams}
\small
\begin{tabular}{ll}
\toprule
\textbf{Hyperparameter} & \textbf{Value} \\
\midrule
\multicolumn{2}{l}{\textit{Graph and feature selection}} \\
\quad Selected genes ($N_g$)        & $700$ \\
\quad Selected miRNAs ($N_m$)       & $100$ \\
\quad Selected TFs ($N_t$)          & $200$ \\
\quad Feature selection             & variance \\
\quad Active edge types             & $E_{GG}, E_{MG}, E_{MM}, E_{TG}, E_{TT}, E_{TM}$ \\
\midrule
\multicolumn{2}{l}{\textit{Architecture}} \\
\quad GNN backbone                  & GCN \\
\quad Pool size                     & $16$ \\
\quad Jumping knowledge             & disabled \\
\quad Parallel branch               & enabled \\
\quad Decoder (reconstruction)      & enabled \\
\quad Dropout                       & $0.2$ \\
\midrule
\multicolumn{2}{l}{\textit{Optimisation}} \\
\quad Optimiser                     & AdamW \\
\quad Learning rate                 & $8 \times 10^{-4}$ \\
\quad Weight decay                  & $1 \times 10^{-4}$ \\
\quad Manual L2 on GNN weights      & $5 \times 10^{-4}$ \\
\quad Batch size                    & $64$ \\
\quad LR schedule                   & cosine annealing \\
\quad Max epochs / patience         & $100$ / $10$ \\
\midrule
\multicolumn{2}{l}{\textit{Class imbalance}} \\
\quad Sampler                       & \texttt{WeightedRandomSampler} \\
\quad Weighting                     & $w_c = K/(\gamma n_c + (1-\gamma)K)$ \\
\quad $\gamma$                      & $0.98$ \\
\quad Cross-entropy weighting       & disabled (handled by sampler) \\
\bottomrule
\end{tabular}
\end{table}

\paragraph{Handling class imbalance.}
The 27-class label distribution is strongly skewed; the largest class
contains roughly ten times more samples than the smallest, and four
classes have fewer than fifty samples each. We mitigate this at the
sampler level rather than through a class-weighted cross-entropy, as
preliminary runs with the latter were unstable on the rare classes.
Each training mini-batch is drawn through a
\texttt{WeightedRandomSampler} that assigns to every sample a class
weight
%
\begin{equation}
    w_c \;=\; \frac{K}{\gamma\, n_c + (1-\gamma)\, K},
    \label{eq:gamma_weights}
\end{equation}
%
where $K$ is the size of the training set, $n_c$ the cardinality of
class $c$, and $\gamma \in [0,1]$ a hyperparameter that interpolates
between random sampling ($\gamma = 0$) and inverse-frequency sampling
($\gamma = 1$). We use $\gamma = 0.98$ in the final configuration, a
value selected after a dedicated sampler ablation
(Section~\ref{sec:results}) that explored the grid
$\gamma \in \{0.90, 0.95, 0.98, 0.99\}$ together with the alternative
inverse-frequency parameterisation
$w_c = (K/n_c)^{1/\alpha}$ for $\alpha \in \{1, 2, 3\}$. The
$\gamma$-parameterisation avoids the divergent weights of pure
inverse-frequency sampling on the smallest classes.

\paragraph{Hyperparameter search.}
We ran a $35$-trial Optuna~\cite{akiba2019optuna} study optimising
validation macro-F1, with a median pruner ($5$ startup trials, $15$
warmup steps) and a hard $10$-hour wall-clock timeout. To make trials
comparable we fixed the random seed and ran each trial for at most
$50$ epochs with early stopping on validation macro-F1 (patience
$7$), then re-trained the selected configuration for $100$ epochs on
the eleven evaluation seeds reported in
Section~\ref{sec:results}. The search covered learning rate, weight
decay, dropout, batch size, the graph-readout pool size, the use of
jumping-knowledge connections, the decoder and parallel-branch on/off
switches, the number of selected genes
$N_g \in \{500, 700, 1000\}$, and the sampler strategy with its
hyperparameters. The categorical choice between GCN and GAT backbones
was kept outside the search space and treated as a separate axis of
the ablation (Section~\ref{sec:results}), so that both architectures
could be compared under their own best configurations. The selected
configuration corresponds to trial 18 of the study and is reproduced
verbatim by the \texttt{config\_final.yml} file in the project
repository.

\paragraph{Evaluation metrics.}
Because the label space is large and imbalanced, we report accuracy,
macro-F1 and weighted-F1, treating macro-F1 as the primary metric.
Per-class precision, recall and F1 are reported in
Section~\ref{sec:results} alongside the aggregate numbers. Statistical
comparisons between models are based on the eleven-seed paired runs:
we use the Friedman test as an omnibus check across all models, and
pairwise Wilcoxon signed-rank tests with Holm correction for multiple
comparisons whenever a $p$-value is associated to a specific model
pair.

\paragraph{Implementation.}
The framework is implemented in PyTorch $2.1.1$ and PyTorch Geometric
$2.4.0$. Experiments were run on a workstation equipped with an
Intel Core i7-10750H CPU, 16~GB RAM and an Nvidia GeForce RTX~2060 GPU
(6~GB VRAM). Full training of a single seed at the final configuration
takes approximately $25$ minutes on this hardware; the complete
ablation suite (five tables, $\sim$\,$450$ runs) takes on the order of
$10$ GPU-days. All code, configuration files and the Makefile-driven
reproduction pipeline are released publicly to allow exact replication.



\begin{comment}

**Cose che ho deliberatamente fatto, con motivazione**

1. **Equazioni numerate per le adiacenze derivate.** Le equazioni $\mathbf{A}_{MM}, \mathbf{A}_{TT}, \mathbf{A}_{TM}$ sono il cuore tecnico del contributo: dimostrano che l'estensione TF non è un'aggiunta ad hoc, ma una derivazione coerente dal blocco TF--gene. Numerarle aiuta i reviewer a riferirsi a loro nel report.

2. **Ho dichiarato il dettaglio dell'overlap genes/TFs ($48$ per $N_g{=}700$).** Era nella tua bozza e si vede dai grafi. Includerlo nei Methods è onesto: lo trovano comunque in figura 4 (graphs), e i reviewer apprezzano che venga discusso prima invece che lasciato a un'osservazione ex-post.

3. **"We do not use edge weights, partly to ... partly because ..."** — qui ho preso una posizione difensiva ragionevole. Le tue analisi mostrano che binarizzare era una limitazione (nella Discussion la dichiari). Anticiparla nei Methods con una giustificazione doppia (compatibilità con baseline + scale eterogenee dei priors) la rende difendibile senza essere apologetica.

4. **"We acknowledge that this is a design choice; collapsing them is a sensible alternative that we did not explore."** sui geni duplicati tra $V_G$ e $V_T$. Hedging onesto. I reviewer di *Bioinformatics* premiano questo tipo di trasparenza.

5. **Patience 20 per early stopping** — non l'ho trovato nei tuoi file, ma è un valore standard e plausibile dato il `num_epochs=100`. **Verifica e correggi se diverso**.

6. **Cosine annealing con warmup di 5 epoche** — ho scelto questa formulazione perché il README dice che lo scheduler è "selectable". **Verifica quale scheduler usa effettivamente `config_final.yml`**: se è costante o un altro, sostituisco con quello.

7. **Layer count del projector = 8 blocchi.** Nella bozza italiana c'è "Lineari 2-8" nella figura; ho interpretato come 8 blocchi, ma potrebbe essere "blocchi da 2 a 8 layer", ovvero un range. **Controlla questo**, perché è il tipo di dettaglio che un reviewer di *Bioinformatics* va a guardare nel codice.

8. **Tempi: 25 min/seed, 10 GPU-day per ablation.** Ho dato due numeri concreti. Sono stime mie basate sulla descrizione dell'hardware. **Se hai numeri reali, sostituiscili**: rende la sezione molto più credibile.

**Cose che NON ho fatto**

- *Non* ho dato la formula esatta del GAT/GCN message passing (sarebbe troppo testo per concetti standard nel target journal). Bastano le citazioni Kipf & Welling 2017 e Veličković et al. 2018, già nella tua bib.
- *Non* ho descritto in dettaglio MoGCN e Li24 come baseline qui: vanno nella sezione Results, dove parli del confronto. Methods descrive il *tuo* metodo, non quelli a confronto.

**Citazioni utilizzate**

Tutte già presenti nella tua `bibliography.bib` (e additions per quelle TODO 1+2+5-6):
- `li_multimodal_2024`, `kipf_semi-supervised_2017`, `velickovic_graph_2018`, `liska_tflink_2022`, `stark_biogrid_2006`, `chen_mirdb_2020`, `xiao_graph_2023`, `akiba2019optuna`.

**Nessuna nuova citazione introdotta** → il file `bibliography_additions.bib` non cambia in questo turno. Resta al contenuto di fine TODO 5-6.

**Verifiche operative che ti chiedo prima di procedere**

| # | Dettaglio | Cosa controllare |
|---|---|---|
| 1 | Patience early stopping | Valore reale (ho messo 20) |
| 2 | LR scheduler effettivo | Cosine + warmup, oppure altro? |
| 3 | Layer projector | 8 fisso, oppure range 2–8? |
| 4 | miRDB score threshold | Ho messo $\geq 80$, è quello che usate? |
| 5 | Numero geni dopo merge ENSEMBL/HGNC | Era `$N_1$` simbolico in TODO 5; se lo sai concretamente, lo metto |
| 6 | Tempo wall-clock per seed | Stima reale al posto della mia? |
\end{comment}
% =====================================================================
% RESULTS — v1, EN
% Da inserire DOPO la sezione Materials and methods.
% Le tabelle riprendono i numeri della bozza; le caption sono riscritte
% in inglese e asciutte. Le tabelle stesse vanno copiate dalla bozza
% (qui riporto solo struttura + caption nuove dove serve).
% =====================================================================


\begin{comment}

Riepilogo di questo turno, con le scelte da segnalarti.

**Cosa ho fatto, e perché — passaggio da tesi a Q1**

1. **Ho ribaltato l'ordine logico.** La tesi presentava prima Methods (incluso metriche, baseline, hyperparameter), poi Results sparsi tra "Performance" e "Ablation study" in capitoli separati e con un cappello discorsivo ("In questo capitolo i risultati vengono discussi alla luce delle domande..."). In un paper Q1 Results apre con il risultato principale e basta. Ho eliminato del tutto i cappelli introduttivi e le frasi tipo "per approfondire il confronto è stata condotta un'analisi". Il lettore Q1 sa già cosa sta leggendo.

2. **Lead-with-the-result in ogni sottosezione.** Ogni paragrafo apre col numero, non con la procedura. "MOGNN-TF reaches 0.7918, the best of all models" invece di "abbiamo confrontato MOGNN-TF con le baseline e i risultati suggeriscono...".

3. **Titoli di sottosezione che sono claim, non etichette.** "The benefit is robust to less supervision", "Connectivity, not edges per se". Questo è lo stile *Bioinformatics*/*Nature Methods*: il titolo dice già cosa hai trovato.

4. **Il non-risultato TF↔miRNA è trattato come finding, non come ripiego.** "The result we want to highlight, though, is a non-difference." Questa è la frase più importante dei Results, e l'ho isolata. Un reviewer Q1 apprezza che un autore sappia che un p-value alto può essere il risultato più interessante.

5. **Ho spostato due tabelle in Supplementary** (per-class S1, sampler per-class S2). In un paper Q1 lo spazio è prezioso: le tabelle giganti per-classe (27 righe × 3 modelli × 3 metriche) e il dettaglio del sampler per-classe vanno in supplementary, citate nel testo. Le due tabelle che restano inline (model comparison, split ratio, gene ablation, TF ablation) sono quelle che sostengono i claim principali.

6. **Ho rivisto la spiegazione GCN>GAT.** La tesi diceva "probabilmente a causa della maggiore sensibilità al rumore topologico delle architetture attention-based". L'ho resa più precisa e difendibile: "GAT has more parameters to fit on the same limited supervision". È un'ipotesi falsificabile, non una congettura vaga.

**Tre cose che ti devo segnalare** 

[Alcuni punti sono da certificare sulla repo] 

1. **Numeri dell'ablation TF su 16 run, non 11.** La `tab:ablation_tf` riporta media su **16 run con sampling random**, mentre la `tab:model_comparison` è su **11 seed**. Questo spiega perché gene+miRNA+TF lì fa 0.7913 mentre il main result fa 0.7918: sono setting diversi (sampler random vs weighted). L'ho dichiarato esplicitamente nella caption ("sixteen runs per row, random sampler") e nel testo. **Verifica che 16 sia il numero giusto** — la bozza dice "16 run" ma vorrei essere sicuro non sia un refuso per 11.

2. **MoGCN F1-macro 0.6473.** Il numero è onesto ma molto basso, e l'ho contestualizzato con il constraint di memoria (2000/2000/300 feature) invece di lasciarlo sembrare un cherry-pick a favore nostro. Questo è esattamente come un reviewer Q1 vuole vedere trattata una baseline che va male: spiegare *perché*, non esultare.

3. **Riferimenti a Supplementary e figure.** Ho citato `Supplementary Table~S1`, `S2`, `Figure~\ref{fig:confusion}` e `Figure~\ref{fig:graphs}c`. La confusion matrix figure (`fig:confusion`) nella tua bozza ha label `fig:CM_mia` — va rinominata coerentemente, oppure cambio io il `\ref`. E le tabelle per-classe vanno effettivamente spostate in un file supplementary se concordi con la scelta.

**Tabelle che NON ho incluso inline** (da decidere): la `tab:sampler_combined` (strategie campionamento, parte alta) e la `tab:overlap_genes_tfs_compact` (overlap geni/TF). La prima la richiamo come S2; la seconda l'ho già menzionata in Methods §Graph (l'overlap di 48). Se preferisci tenerne una inline, dimmelo.

**Una domanda di scoping per la qualità Q1**

Il punto più debole dei Results, per un reviewer Q1, è che il vantaggio sulla baseline più vicina (Li24-GCN, +0.0215) è reale e significativo ma modesto, e che il finding TF↔miRNA è un *non-risultato* (per quanto interessante). Due opzioni su come inquadrarlo nella Discussion:
- (a) Tenere il framing attuale, onesto: "small but significant improvement + a substitutability finding with clinical value".
- (b) Spingere di più sul finding di sostituibilità come contributo principale, retrocedendo il +0.0215 a contorno.

Personalmente raccomando (a): è più difendibile e meno attaccabile in review. Ma è una decisione che incide su Discussion e Abstract, quindi vale la pena allinearci prima di scrivere la Discussion.

**Nessuna citazione nuova** → `bibliography_additions.bib` invariato.

Procediamo con **Discussion + Conclusions (TODO 10+11)**? È il punto dove le due domande inquadranti e il flusso causale fanno la differenza tra una tesi e un paper.
\end{comment}



\section{Results}
\label{sec:results}

We evaluate MOGNN-TF on the 27-class iCluster pan-cancer task against
two published baselines, then dissect the model along four axes: the
amount of input information (number of genes), the sampling strategy,
the composition of the prior-knowledge graph, and the amount of
supervision (training fraction). Unless stated otherwise, all numbers
are mean $\pm$ standard deviation over eleven seeds on the held-out
test set, and macro-F1 is the primary metric. Significance is assessed
with a Friedman omnibus test followed by Holm-corrected pairwise
Wilcoxon signed-rank tests at $\alpha = 0.05$.

\subsection{MOGNN-TF outperforms graph-based baselines}
\label{ssec:res_baseline}

MOGNN-TF with a GCN backbone reaches a test macro-F1 of
$0.7918 \pm 0.0156$, the best of all models compared
(Table~\ref{tab:model_comparison}). It improves over the prior-knowledge
baseline of Li and Nabavi~\cite{li_multimodal_2024} by $+0.0215$ on its
GCN variant ($0.7703$) and by $+0.0361$ on its GAT variant ($0.7557$),
and by a large margin over the patient-similarity baseline
MoGCN~\cite{li_mogcn_2022} ($0.6473$). The Friedman test rejects the
null hypothesis of equal performance across the five models
($\chi^2 = 24.79$, $p = 1.5\times10^{-4}$), and the post-hoc Wilcoxon
analysis confirms that MOGNN-TF (GCN) is significantly better than
Li24 (GAT) ($p_{\text{Holm}} = 0.0029$), Li24 (GCN)
($p_{\text{Holm}} = 0.0059$), MoGCN ($p_{\text{Holm}} = 0.0029$) and
its own GAT counterpart ($p_{\text{Holm}} = 0.0205$). The gap over
MoGCN is the largest, and is partly explained by a constraint of that
method rather than of ours: building its patient-similarity network is
memory-bound, so MoGCN could not be run on the full feature set and was
restricted to a variance-pruned input of 2{,}000 genes, 2{,}000 TFs
and 300 miRNAs.

\begin{table}[!t]
\centering
\caption{Test-set performance on the 27-class pan-cancer task, mean
$\pm$ s.d.\ over eleven seeds. $p$-values are Holm-corrected Wilcoxon
signed-rank tests against MOGNN-TF (GCN) on macro-F1. Best value in
bold.}
\label{tab:model_comparison}
\begin{adjustbox}{max width=\linewidth}
\begin{tabular}{lccccc}
\toprule
\textbf{Model} & \textbf{Accuracy} & \textbf{F1 Macro} & \textbf{F1 Weighted} & \textbf{$p$} & \textbf{$p_{\text{Holm}}$} \\
\midrule
Li24 (GAT)~\cite{li_multimodal_2024} & $0.8077 \pm 0.0155$ & $0.7557 \pm 0.0169$ & $0.8024 \pm 0.0152$ & $4.9\times10^{-4}$ & $0.0029$ \\
Li24 (GCN)~\cite{li_multimodal_2024} & $0.8220 \pm 0.0119$ & $0.7703 \pm 0.0191$ & $0.8169 \pm 0.0126$ & $1.5\times10^{-3}$ & $0.0059$ \\
MoGCN~\cite{li_mogcn_2022}           & $0.7843 \pm 0.0077$ & $0.6473 \pm 0.0084$ & $0.7552 \pm 0.0080$ & $4.9\times10^{-4}$ & $0.0029$ \\
MOGNN-TF (GAT)                       & $0.8200 \pm 0.0113$ & $0.7824 \pm 0.0155$ & $0.8188 \pm 0.0110$ & $6.8\times10^{-3}$ & $0.0205$ \\
MOGNN-TF (GCN)                       & $\mathbf{0.8285 \pm 0.0134}$ & $\mathbf{0.7918 \pm 0.0156}$ & $\mathbf{0.8264 \pm 0.0121}$ & --- & --- \\
\bottomrule
\end{tabular}
\end{adjustbox}
\end{table}

The GCN backbone is consistently ahead of GAT, here and in every
ablation below. We read this as a sign that, on a graph this sparse and
noisy after feature selection, the smoother neighbourhood averaging of
GCN is more stable than the learned attention of GAT, which has more
parameters to fit on the same limited supervision. Per-class precision,
recall and F1 are reported in Supplementary
Table~S1, and the normalised confusion matrices of MOGNN-TF and the
Li24 baseline are shown in Supplementary Figure~S1; most of the
residual error concentrates on a small number of low-support,
biologically adjacent clusters.

\subsection{The benefit is robust to less supervision}
\label{ssec:res_split}

To probe behaviour under the data scarcity typical of clinical cohorts,
we retrained the model with training fractions of $80\%$ down to $50\%$,
holding the validation set at $10\%$ and assigning the remainder to test
(Table~\ref{tab:splits_comparison}). Performance degrades gracefully:
the GCN macro-F1 moves from $0.7978$ at $80\%$ to $0.7779$ at $50\%$,
and the GCN--GAT ordering is preserved at every fraction. The
degradation is not statistically resolvable in the upper range — the
difference between the $80\%$ and $60\%$ GCN models is not significant
($p_{\text{Holm}} = 0.83$) — which suggests that the biologically
constrained graph acts as a regulariser strong enough that the model
does not need the full training set to reach a competitive
representation.

\begin{table}[!t]
\centering
\caption{Effect of the training fraction on the proposed model
(mean $\pm$ s.d., eleven seeds). The validation set is fixed at $10\%$.}
\label{tab:splits_comparison}
\small
\begin{tabular}{llccc}
\toprule
\textbf{Train \%} & \textbf{Model} & \textbf{Accuracy} & \textbf{F1 Macro} & \textbf{F1 Weighted} \\
\midrule
\multirow{2}{*}{$80\%$} & GAT & $0.8287 \pm 0.0094$ & $0.7891 \pm 0.0205$ & $0.8255 \pm 0.0108$ \\
                        & GCN & $0.8355 \pm 0.0123$ & $0.7978 \pm 0.0205$ & $0.8342 \pm 0.0119$ \\
\midrule
\multirow{2}{*}{$70\%$} & GAT & $0.8181 \pm 0.0164$ & $0.7798 \pm 0.0216$ & $0.8153 \pm 0.0171$ \\
                        & GCN & $0.8245 \pm 0.0132$ & $0.7865 \pm 0.0155$ & $0.8224 \pm 0.0119$ \\
\midrule
\multirow{2}{*}{$60\%$} & GAT & $0.8156 \pm 0.0141$ & $0.7756 \pm 0.0131$ & $0.8133 \pm 0.0139$ \\
                        & GCN & $0.8292 \pm 0.0144$ & $0.7971 \pm 0.0162$ & $0.8280 \pm 0.0142$ \\
\midrule
\multirow{2}{*}{$50\%$} & GAT & $0.8106 \pm 0.0178$ & $0.7735 \pm 0.0187$ & $0.8091 \pm 0.0162$ \\
                        & GCN & $0.8169 \pm 0.0151$ & $0.7779 \pm 0.0178$ & $0.8154 \pm 0.0150$ \\
\bottomrule
\end{tabular}
\end{table}

\subsection{Ablations}
\label{ssec:res_ablation}

\paragraph{Number of input genes.}
Macro-F1 is highest at $700$ genes for both backbones
(Supplementary Table~S3), with GCN (700) reaching
$0.7918 \pm 0.0156$. Moving to $500$ or $1{,}000$ genes gives
statistically indistinguishable results (GCN-500 and GCN-1000 are not
significantly different from GCN-700 after correction), while $300$
genes is clearly worse ($p_{\text{Holm}} \leq 0.017$ for both
backbones). The plateau between $500$ and $1{,}000$ genes is
informative: beyond a few hundred high-variance genes, adding nodes
does not add discriminative signal, and the gain we report does not
come from simply enlarging the graph.

\paragraph{Sampling strategy.}
A Friedman test across eight sampling strategies (seven seeds) on
validation macro-F1 finds no globally significant difference
($\chi^2 = 9.15$, $p = 0.24$): no single strategy dominates on the
aggregate metric. The weighted sampler nevertheless helps where it is
designed to. A per-class comparison between $\gamma = 0.98$ and random
sampling (Supplementary Table~S2) shows consistent recall gains on the
lowest-support clusters — for example C1 (10 samples,
$\Delta\text{recall} = +0.114$) and C2 (30 samples,
$\Delta\text{recall} = +0.095$) — at the cost of small, mostly
non-significant losses on well-populated classes. We keep
$\gamma = 0.98$ in the final configuration on this basis: the average
per-class recall gain is positive ($+0.034$ on validation, $+0.023$ on
test) even where the aggregate metric is flat.

\paragraph{Prior-knowledge composition.}
The most informative ablation concerns which biological priors enter
the graph (Table~\ref{tab:ablation_tf}; sixteen runs per row, random
sampler, to isolate the effect of the graph from that of the sampler).
Adding the full TF layer on top of the gene+miRNA graph gives a small
improvement in GCN macro-F1, from $0.7843$ to $0.7913$, and the
overall Friedman test across configurations is significant
($\chi^2 = 31.22$, $p = 2.7\times10^{-4}$). The result we want to
highlight, though, is a non-difference. Replacing miRNAs with TFs
entirely — a gene+TF graph with no miRNA layer — yields
$0.7866 \pm 0.0148$ for GCN, statistically indistinguishable from the
gene+miRNA graph at $0.7843 \pm 0.0138$
($\Delta = +0.0023$, $p_{\text{Holm}} = 1.00$); the same holds for GAT
($\Delta = -0.0032$, $p_{\text{Holm}} = 1.00$). In other words, the
transcriptional layer carries enough of the regulatory signal that it
can stand in for the post-transcriptional one without a measurable
loss. Since TF expression is already contained in the transcriptomic
readout, this substitution adds no new assay; it makes explicit, as
graph topology, information that was latent in gene expression. This
matters for the common clinical situation where miRNA profiling is
absent (block missingness): a gene+TF graph recovers comparable
performance from RNA-seq alone.

\begin{table}[!t]
\centering
\caption{Prior-knowledge ablation (mean $\pm$ s.d., sixteen runs per
row, random sampler). Replacing miRNAs with TFs (gene+TF) matches the
gene+miRNA graph; the differences are not significant after Holm
correction. Best per column in bold.}
\label{tab:ablation_tf}
% \resizebox{\linewidth}{!}{%
\begin{tabular}{llcccc}
\toprule
\textbf{Configuration} & \textbf{Model} & \textbf{Genes} & \textbf{TF} & \textbf{miRNA} & \textbf{F1 Macro} \\
\midrule
\multirow{2}{*}{Gene+miRNA}    & GCN & \checkmark & --        & \checkmark & $0.7843 \pm 0.0138$ \\
                               & GAT & \checkmark & --        & \checkmark & $0.7776 \pm 0.0159$ \\
\midrule
\multirow{2}{*}{Gene+miRNA+TF} & GCN & \checkmark & \checkmark & \checkmark & $\mathbf{0.7913 \pm 0.0112}$ \\
                               & GAT & \checkmark & \checkmark & \checkmark & $0.7777 \pm 0.0131$ \\
\midrule
\multirow{2}{*}{Gene+TF}       & GCN & \checkmark & \checkmark & --        & $0.7866 \pm 0.0148$ \\
                               & GAT & \checkmark & \checkmark & --        & $0.7744 \pm 0.0145$ \\
\midrule
\multirow{2}{*}{Gene only}     & GCN & \checkmark & --        & --        & $0.7856 \pm 0.0143$ \\
                               & GAT & \checkmark & --        & --        & $0.7776 \pm 0.0156$ \\
\bottomrule
\end{tabular}%
% }
\end{table}

\paragraph{Connectivity, not edges per se.}
A complementary observation ties the ablations together. After feature
selection the graph is sparse and contains many isolated nodes
(Figure~\ref{fig:graphs}): for those nodes message passing is inactive
and the representation collapses to the input features. Adding the TF
layer increases the connectivity of the surviving subgraph
(Figure~\ref{fig:graphs}c), which is the most plausible reason a small
but consistent gain appears when TFs are included. We tested several
connectivity-aware feature-selection schemes (signal-propagation
ranking, connectivity refinement) to push this further, but none
improved over plain variance selection; favouring well-connected nodes
tended to select structurally central but less discriminative genes,
trading input signal for topology. We therefore report variance
selection throughout, and treat the connectivity effect as an
explanation of the TF benefit rather than a lever to optimise.



\begin{comment}

Leggo prima la Discussion della tesi, così ragiono sul testo reale e non in astratto.Ho letto tutta la Discussion della tesi. Ora rispondo prima al nodo di scoping, perché hai ragione e voglio spiegarti *perché* hai ragione — e anche dove serve una sfumatura.

## Sul nodo di framing: hai ragione, ma con una precisazione importante

La tua posizione è corretta: non devo *scegliere* tra (a) e (b) come se fossero alternative. Sono due contributi reali e distinti, e un paper Q1 forte li rivendica entrambi. La mia impostazione iniziale era troppo difensiva — pensavo "in review" invece di "valore del lavoro". Difendere il valore del lavoro è la priorità, e il lavoro ha *due* risultati, non uno.

La precisazione che però conta, e che non è una ritirata sulla mia raccomandazione ma un raffinamento, è questa: **i due risultati vivono su piani diversi, e proprio per questo non competono — si rinforzano.** Vanno rivendicati entrambi, ma inquadrati per quello che ciascuno è:

- **Risultato 1 — il miglioramento di performance (+0.0215 su Li24-GCN).** È un risultato *comparativo* e *statisticamente significativo*. Il punto di forza non è la grandezza del delta (che è modesta), ma il fatto che sia significativo dopo correzione di Holm su 11 seed, **e che sia ottenuto sul setting più difficile della letteratura** (27 classi pan-cancer, non BRCA-5-classi). Il framing corretto non è "abbiamo migliorato di poco", ma "abbiamo migliorato in modo statisticamente solido su un task dove la maggior parte dei metodi non si misura nemmeno".

- **Risultato 2 — la sostituibilità TF↔miRNA.** È un risultato *meccanicistico* e *clinicamente azionabile*. È un non-rifiuto dell'ipotesi nulla, sì, ma è un non-rifiuto *informativo*: dice che il livello trascrizionale porta abbastanza segnale regolatorio da rimpiazzare quello post-trascrizionale. Questo non si misura in punti di F1 — il suo valore è in cosa *implica* (RNA-seq da solo basta quando manca il miRNA profiling).

Il motivo per cui i due si rinforzano anziché competere: il Risultato 2 **spiega** il Risultato 1. Il TF layer migliora la performance (R1) *perché* densifica il grafo e consolida segnale regolatorio latente (la connettività), e quella stessa proprietà è ciò che rende i TF intercambiabili coi miRNA (R2). C'è un terzo elemento nella tua Discussion che lega i due — l'osservazione sulla connettività del grafo dopo feature selection — ed è il ponte concettuale. È il filo che trasforma "due risultati separati" in "una tesi coerente".

Quindi la mia raccomandazione operativa, rivista alla luce della tua giusta obiezione: **uso entrambi, do a ciascuno il piano che gli compete, e li lego con la connettività come meccanismo unificante.** Concretamente, le due domande inquadranti della Discussion (come prescrive lo style guide) diventano:

1. *Does an explicit transcription-regulation layer improve pan-cancer subtype classification, and is the gain attributable to biology or to graph topology?* → risponde R1 + connettività.
2. *Can the transcriptional layer stand in for the post-transcriptional one, and what does that imply for settings where miRNA data are missing?* → risponde R2.

Questo prende **tutto** il valore del lavoro — performance, sostituibilità, robustezza allo split, ruolo della connettività — senza buttare via nulla e senza sovravendere niente.

Prima di scrivere, una sola conferma su un punto che cambia il testo in modo sostanziale:
\end{comment}


\begin{comment}
Riepilogo, con le scelte che voglio dichiararti esplicitamente.

**Come ho gestito il tuo nodo di scoping.** Entrambi i contributi sono rivendicati, ciascuno sul suo piano. R1 (performance) è inquadrato come "small but statistically robust, on a task most methods avoid" — il valore è la solidità statistica e la difficoltà del setting, non il delta. R2 (sostituibilità) come "more than a technical curiosity", con l'implicazione clinica esplicita. E la connettività del grafo è il ponte: la stessa proprietà del TF layer spiega entrambi. Niente è stato buttato via.

**Cosa ho fatto, e perché — passaggio da tesi a Q1.**
1. **Eliminato l'elenco puntato "Discussion of activities".** La tesi enumerava 5 "attività" in un `enumerate`. In un paper Q1 i contributi stanno nell'Introduction; la Discussion *discute*, non ri-elenca. Ho convertito il contenuto utile di quei 5 punti in prosa argomentativa dentro i paragrafi tematici.
2. **Due domande inquadranti in apertura**, come prescrive lo style guide, e ogni paragrafo risponde a una.
3. **Titoli di paragrafo che sono affermazioni**, non etichette ("The gain is real, modest, and topological in origin", "Transcription factors can substitute for miRNAs", "Where the model still fails, and why").
4. **Limitazioni integrate come paragrafo della Discussion**, non sezione separata — più asciutto per Q1. Ho tenuto tutti e cinque i limiti reali della tesi (prior generici, edge binari, feature-selection/connettività, dataset piccolo + no validazione esterna, costo computazionale) ma in prosa densa invece che a blocchi.
5. **Conclusion ridotta a ~180 parole, un paragrafo**, con future work in una sola frase finale invece della sotto-sezione "Sviluppi futuri" estesa. Ho eliminato il ri-elenco dei numeri (0.8285/0.7918/0.8264) perché sono già nei Results: una Conclusion Q1 sintetizza il *significato*, non ripete le cifre.

**Il commento `%` che hai chiesto.** È sopra il paragrafo "Attributing the gain to the TF layer" (il confronto con baseline-modified). Spiega: perché tenerlo (è l'ablation più onesta del contributo TF a livello full-pipeline, anticipa la domanda ovvia del reviewer), qual è il rischio (mostra che a quel livello il delta TF non è significativo), e come rimuoverlo se serve (tagliare il paragrafo e relegare il punto a una frase nei Limiti — R2 non ne dipende).

**Tre cose da segnalarti.**
1. **Citazione `li_cancer_2021`** — usata nel paragrafo sulle classi misclassificate, ripresa dalla tesi. Verifica che sia nella tua `.bib`; se manca te la aggiungo all'addition file (mi serve sapere a quale paper si riferisce — probabilmente un GNN su pan-cancer subtype 2021).
2. **Il numero "diagonal mass C1 da 0.427 a 0.645"** viene dalla tesi (confronto sampler weighted vs random sulla GCN). L'ho riportato come esempio concreto; verifica che 0.427→0.645 sia il confronto giusto (nella tesi è il passaggio da Li24-GCN recall 0.427 a MOGNN-TF-GCN recall 0.645 su C1, coerente con la tabella S1).
3. **Nessuna citazione nuova introdotta** oltre a quelle già discusse → `bibliography_additions.bib` invariato in questo turno (a meno che `li_cancer_2021` non manchi).

A questo punto abbiamo in piedi tutte le sezioni testuali: Abstract, Introduction, Related Work, Methods (completo), Results, Discussion, Conclusions, più il Supplementary. Vuoi che adesso facccia un **check di coerenza globale** su tutto il manoscritto (allineamento numeri, riferimenti, claim tra sezioni), oppure preferisci affrontare il **pass di umanizzazione finale** (TODO 13) o la verifica della bibliografia completa (TODO 3)?
\end{comment}




% =====================================================================
% DISCUSSION — v1, EN  (TODO 10)
% Sostituisce l'intera \section{Discussion} della tesi (cappello +
% Performance + Feature selection + Discussion of activities a enumerate
% + Limitazioni + Contributo originale).
%
% Struttura Q1: due domande inquadranti, flusso causale, NO elenco
% puntato di "attività". I contributi sono già nell'Introduction; qui
% si DISCUTE, non si ri-enumera.
% =====================================================================
\section{Discussion}
\label{sec:discussion}

Two questions organise this section. First, does an explicit
transcription-regulation layer improve pan-cancer subtype
classification, and is any gain attributable to the biology of the
prior or to the topology it induces? Second, can the transcriptional
layer stand in for the post-transcriptional one, and what would that
mean for settings where miRNA profiles are missing? The two are
connected: as we argue below, the same property of the TF layer — that
it reconnects a graph fragmented by feature selection — accounts both
for the performance gain and for the substitutability we observe.

\paragraph{The gain is real, modest, and topological in origin.}
MOGNN-TF (GCN) improves over the published Li24 baselines by a margin
that is small in absolute terms ($+0.0215$ macro-F1 over Li24-GCN) but
statistically robust across eleven seeds, on a 27-class taxonomy that
most prior-knowledge GNNs do not attempt. We read the size of the gain
as informative rather than disappointing. The ablations
(Section~\ref{sec:results}) show that adding nodes beyond $\sim$700
high-variance genes does not help, and that replacing the sampler does
not move the aggregate metric; what does change the picture is the
connectivity of the graph that survives feature selection. Variance-based
selection, applied per node type, fragments the prior network and leaves
many isolated nodes for which message passing is inactive
(Figure~\ref{fig:graphs}). Introducing the TF layer reconnects a portion
of this subgraph, and the performance improvement tracks that
reconnection. The implication is that the benefit of biological prior
knowledge in this setting is not automatic — it does not follow from the
mere presence of curated edges — but is mediated by how much usable
structure survives dimensionality reduction. This also explains a
negative result of our own: connectivity-aware selection schemes
(signal-propagation ranking, connectivity refinement) did not beat plain
variance selection, because favouring well-connected nodes traded away
discriminative input signal for topology. Edges help only when they
connect nodes that are themselves informative.

% --- NOTE (review) -------------------------------------------------
% The paragraph below compares the final model against the "baseline
% modified" (our framework WITHOUT the TF layer, gene+miRNA only).
% Rationale for keeping it: it is the most honest possible ablation of
% the TF contribution at the FULL-pipeline level, and pre-empts the
% obvious reviewer question "is the gain really the TF layer or the
% other implementation changes?". Risk: it shows the TF gain is NOT
% significant at this level (only at the prior-composition ablation
% level, Table tab:ablation_tf), which a hostile reviewer may use to
% downplay R1. If the paper needs to read more assertively, this
% paragraph can be cut and the point relegated to a single sentence in
% Limitations; the substitutability result (R2) does not depend on it.
% -------------------------------------------------------------------
\paragraph{Attributing the gain to the TF layer.}
A stricter test isolates the TF layer inside our own pipeline. When the
final model is compared against the same framework restricted to
gene--miRNA edges (the ``baseline modified'' configuration), the
macro-F1 advantage is no longer significant after Holm correction.
We do not regard this as undercutting the contribution: at the level
of the prior-knowledge ablation (Table~\ref{tab:ablation_tf}), where the
sampler and other factors are held fixed, the full TF graph is the best
configuration, and the per-class analysis shows where the gain
concentrates. But it does set the honest boundary of the claim — the
TF layer yields a consistent, not a large, improvement, and its main
practical value lies in the substitutability result rather than in
raw macro-F1.

\paragraph{Transcription factors can substitute for miRNAs.}
The clearest finding of the ablation is a non-difference. A graph built
from genes and TFs alone, with no miRNA layer, matches the gene+miRNA
graph on macro-F1, with no statistically significant difference for
either backbone (GCN: $\Delta = +0.0023$, $p_{\text{Holm}} = 1.00$;
GAT: $\Delta = -0.0032$, $p_{\text{Holm}} = 1.00$). This is more than a
technical curiosity. TF expression is already measured in any RNA-seq
assay, so encoding the transcriptional layer adds no new experiment; it
makes explicit, as graph structure, regulatory information that is
otherwise latent in gene expression. The practical reading is direct:
in the common clinical situation where miRNA profiling is unavailable or
incomplete — a form of block missingness — a gene+TF graph recovers
comparable performance from transcriptomic data alone. We are cautious
about over-generalising from a single taxonomy and a single cohort, but
the result suggests that the post-transcriptional layer, often treated
as essential in multi-omics integration, may be partly redundant once
transcriptional regulation is modelled explicitly.

\paragraph{Where the model still fails, and why.}
The residual error is concentrated, and concentrated where biology
predicts. The hardest classes are C7 and C13 — mixed, copy-number-driven
clusters that are intrinsically heterogeneous — and C17, a BRCA subtype
adjacent to other breast clusters in molecular space; these are among
the most frequently misclassified subtypes for GNNs on this taxonomy in
general~\cite{xiao_graph_2023, li_cancer_2021}. The weighted sampler
($\gamma = 0.98$) improves recall on several low-support classes (for
example C1, from $0.427$ to $0.645$ diagonal mass) without resolving the
mixed-aneuploidy classes, which is consistent with our reading of the
failure mode: for those subtypes, variance-based selection tends to
discard the very genes that would discriminate them, so no amount of
rebalancing at the sampler level can recover a signal that is not in the
input. The robustness experiment points the same way — performance is
stable from $80\%$ down to $60\%$ training data
($p_{\text{Holm}} = 0.83$ between the two), which we attribute to the
biological graph acting as a topological regulariser that reduces the
model's dependence on large labelled sets.

\paragraph{Limitations.}
% Versione 1 (prima del TODO 5)
% Several constraints bound these conclusions. The prior-knowledge
% networks are generic rather than cancer- or context-specific: BioGRID,
% miRDB and TFLink aggregate interactions across tissues and conditions,
% may contain false positives, and certainly miss unannotated edges,
% introducing topological noise. We also treat all edges as binary and
% unweighted, which flattens interaction strength that could otherwise
% guide message passing. The interaction between feature selection and
% graph integrity, discussed above, is itself a limitation: an a-priori
% variance cut fragments the network, and we did not find a selection
% scheme that preserves both discriminative signal and connectivity.
% On the data side, even the TCGA Pan-Cancer Atlas yields only $\sim$9{,}000
% tumours with all three omic layers, the rarest subtypes have very few
% instances, and — most importantly — we evaluate only on internal
% stratified splits of a single consortium. Without external validation on
% an independent cohort we cannot rule out that the learned
% representations have absorbed TCGA-specific batch or platform effects,
% so cross-cohort generalisation remains open. Finally, training on
% heterogeneous multi-level graphs is computationally heavy, and the cost
% grows quickly with the number of nodes, particularly for the
% attention-based variant.

% =============================

% Versione 2 (dopo TODO 5, con roadmap più generica e meno leave-one-TSS-out):
Several constraints bound these conclusions. The prior-knowledge
networks are generic rather than cancer- or context-specific: BioGRID,
miRDB and TFLink aggregate interactions across tissues and conditions,
may contain false positives, and certainly miss unannotated edges,
introducing topological noise. We also treat all edges as binary and
unweighted, which flattens interaction strength that could otherwise
guide message passing. The interaction between feature selection and
graph integrity, discussed above, is itself a limitation: an a-priori
variance cut fragments the network, and we did not find a selection
scheme that preserves both discriminative signal and connectivity.
The most consequential limitation is on the data side. Even the TCGA
Pan-Cancer Atlas yields only $9{,}666$ tumours with all three omic layers
and a matched iCluster label, the rarest subtypes have very few instances,
and we evaluate exclusively on internal stratified splits of a single
consortium. Our evidence for internal validity is deliberately
conservative: performance is reported as macro-F1 $0.7918 \pm 0.0156$
across eleven random seeds, with comparisons established through a
Friedman test followed by Holm-corrected Wilcoxon post-hoc tests, and the
result is preserved across five additional independent stratified data
partitions (macro-F1 $0.780 \pm 0.039$), confirming it is not an artefact
of the particular split.
% Esperimento 1 (stabilità su split multipli): FATTO e confermato.
% Script di riproducibilità: scripts/run_split_stability.py
% Analisi completa: results/pancan/split_stability/ANALYSIS.md
Internal resampling, however, cannot substitute for an independent cohort,
and we cannot rule out that the learned representations have partly
absorbed TCGA-specific batch or platform effects. This constraint is
shared by the graph-based multi-omics models we build on, which are
likewise developed and evaluated within TCGA
\cite{li_mogcn_2022,wang_mogonet_2021,kesimoglu_supreme_2022,li_multimodal_2024};
the studies that do validate externally do so on single-cancer tasks for
which a label-compatible independent cohort exists \cite{tanvir_mogat_2024},
an option the 27-class pan-cancer iCluster taxonomy does not afford, since
no external resource carries those integrated subtype labels. Cross-cohort
generalisation, and robustness to site- and platform-level shift, therefore
remain open and are the natural next step. Finally, training on
heterogeneous multi-level graphs is computationally heavy, and the cost
grows quickly with the number of nodes, particularly for the
attention-based variant.


% =====================================================================
% CONCLUSIONS — v1, EN  (TODO 11)
% Sostituisce \section{Conclusions} + \subsection*{Sviluppi futuri}.
% Un solo paragrafo compatto (~180 parole), future work integrato in
% una frase. NO ri-elenco di numeri già nei Results.
% =====================================================================
\section{Conclusions}
\label{sec:conclusions}

We presented MOGNN-TF, a heterogeneous graph neural network that adds an
explicit transcription-regulation layer to a multi-omics prior-knowledge
graph and evaluates it, as a single integrated structure, on the full
27-class iCluster pan-cancer taxonomy. On this hard, imbalanced setting
the model improves over published graph-based baselines by a small but
statistically robust margin, and an ablation across prior-knowledge
sources shows that transcription factors can replace miRNAs in the graph
without a measurable loss in macro-F1 — a substitution that requires no
additional assay, since TF expression is already present in RNA-seq.
Taken together, the two results indicate that the value of biological
prior knowledge in this task is mediated by the connectivity of the
graph that survives feature selection, rather than by the choice of edge
source alone. The most promising extensions follow directly: relational
or heterogeneous GNN architectures that preserve edge semantics rather
than aggregating them, richer cancer-specific priors such as pathway
networks, connectivity-aware feature selection, and — the prerequisite
for clinical relevance — validation on independent external cohorts.

\section*{Declaration of Competing Interest}
The authors declare that they have no known competing financial interests or personal relationships that could have appeared to influence the work reported in this paper.

% \section*{Ethics statement}
% This study used de-identified, publicly available data from the
% NCI Genomic Data Commons (TCGA) and from the Therapeutically Applicable
% Research to Generate Effective Treatments (TARGET) initiative.
% Institutional Review Board approval was not required.

\section*{Data and Code Availability}
The MOGNN-TF source code, configuration files and the scripts that
reproduce the preprocessing, training and evaluation pipeline are
released under an open-source licence at
\url{https://github.com/johndef64/MOGNN-TF}.
All input data come from public repositories. Gene expression
(RNA-seq), gene-level copy number and miRNA profiles for the TCGA
Pan-Cancer cohort were retrieved from UCSC Xena
(\url{https://xenabrowser.net/}). The prior-knowledge networks are
publicly available: protein--protein interactions from BioGRID
(\url{https://thebiogrid.org/}), miRNA--target predictions from miRDB
(\url{https://mirdb.org/}), and transcription-factor--target regulatory
edges from TFLink (\url{https://tflink.net/}).


\begin{footnotesize}

    \bibliographystyle{elsarticle-num}

    \bibliography{bibliography.bib}

\end{footnotesize}


\end{document}